% ============================================================================
% SPANISCH LANGENTWURF - Sequenz 07, Block b: Números 21-100
% ============================================================================
% Fachfarbe: #c41e3a (Karminrot)
% Tier 1 (vollständig)
% ============================================================================

\documentclass[11pt,a4paper]{article}

% ============================================================================
% PACKAGES
% ============================================================================
\usepackage[utf8]{inputenc}
\usepackage[T1]{fontenc}
\usepackage[ngerman]{babel}
\usepackage{geometry}
\usepackage{fancyhdr}
\usepackage{lastpage}
\usepackage{xcolor}
\usepackage{tcolorbox}
\usepackage{enumitem}
\usepackage{tabularx}
\usepackage{longtable}
\usepackage{array}
\usepackage{booktabs}
\usepackage{hyperref}
\usepackage{ifthen}
\usepackage{amssymb}

% ============================================================================
% KONFIGURATION
% ============================================================================

% --- TIER (immer 1 für Spanisch) ---
\newcommand{\plantier}{1}

% --- FACH ---
\newcommand{\fachid}{spanisch}
\newcommand{\fach}{Spanisch}
\newcommand{\fachfarbe}{c41e3a}

% ============================================================================
% VARIABLEN
% ============================================================================

% --- Metadaten ---
\newcommand{\jahrgangsstufe}{7}
\newcommand{\schulart}{Gesamtschule}
\newcommand{\stundenthema}{Números 21-100 y precios}
\newcommand{\reihenthema}{De compras}
\newcommand{\stundennummer}{7b}
\newcommand{\reihenstunden}{4}
\newcommand{\unterrichtsdatum}{XX.XX.XXXX}
\newcommand{\unterrichtszeit}{XX:XX -- XX:XX}
\newcommand{\raum}{XXX}

% --- Lehrkraft ---
\newcommand{\lehrkraft}{N.N.}
\newcommand{\ausbildungsstand}{Lehrkraft}

% --- Lerngruppe ---
\newcommand{\lerngruppengroesse}{24}
\newcommand{\maennlich}{12}
\newcommand{\weiblich}{12}
\newcommand{\divers}{0}

% --- Kompetenzen/Niveau ---
\newcommand{\gerniveau}{A1}
\newcommand{\lehrwerk}{Qué pasa 1}
\newcommand{\lehrwerkseite}{S. 48-49}

% ============================================================================
% LAYOUT
% ============================================================================
\geometry{left=2.5cm, right=2.5cm, top=2.5cm, bottom=2.5cm}
\definecolor{fach}{HTML}{\fachfarbe}
\definecolor{gray}{HTML}{666666}
\definecolor{lightgray}{HTML}{f5f5f5}
\definecolor{successgreen}{HTML}{2a7a4b}
\definecolor{warningorange}{HTML}{b35c00}

\tcbuselibrary{skins,breakable}

% Farbige Boxen
\newtcolorbox{infobox}[1][]{
    colback=lightgray,
    colframe=fach,
    fonttitle=\bfseries,
    title=#1,
    breakable
}

\newtcolorbox{scorebox}{
    colback=white,
    colframe=fach,
    fonttitle=\bfseries,
    title=Didaktik-Score,
    breakable
}

% Kopf-/Fußzeile
\pagestyle{fancy}
\fancyhf{}
\fancyhead[L]{\small\textcolor{gray}{\fach{} -- Jg. \jahrgangsstufe{} -- \stundenthema}}
\fancyhead[R]{\small\textcolor{gray}{Langentwurf Tier 1}}
\fancyfoot[C]{\small\thepage{} / \pageref{LastPage}}
\renewcommand{\headrulewidth}{0.4pt}
\renewcommand{\footrulewidth}{0pt}

% ============================================================================
% DOKUMENT
% ============================================================================
\begin{document}

% ============================================================================
% DECKBLATT
% ============================================================================
\begin{titlepage}
    \centering
    \vspace*{2cm}

    {\large\textcolor{gray}{\schulart}}\\[2cm]

    {\Huge\textcolor{fach}{\textbf{Langentwurf}}}\\[0.5cm]
    {\large Tier 1 -- Vollständig}\\[2cm]

    {\LARGE\textbf{\stundenthema}}\\[0.5cm]
    {\large im Rahmen der Unterrichtsreihe}\\[0.3cm]
    {\Large\textit{\reihenthema}}\\[2cm]

    \begin{tabular}{rl}
        \textbf{Fach:} & \fach{} (\gerniveau) \\[0.3cm]
        \textbf{Klasse:} & \jahrgangsstufe \\[0.3cm]
        \textbf{Lehrwerk:} & \lehrwerk, \lehrwerkseite \\[0.3cm]
        \textbf{Datum:} & \underline{\hspace{3cm}} \\[0.3cm]
        \textbf{Zeit:} & \underline{\hspace{3cm}} (Doppelstunde) \\[0.3cm]
        \textbf{Raum:} & \underline{\hspace{3cm}} \\[0.3cm]
    \end{tabular}

    \vfill

    {\large Lehrkraft: \lehrkraft}\\[0.3cm]
    {\normalsize\textcolor{gray}{\ausbildungsstand}}\\[1cm]

\end{titlepage}

% ============================================================================
% INHALTSVERZEICHNIS
% ============================================================================
\tableofcontents
\newpage

% ============================================================================
% 1. BEDINGUNGSANALYSE
% ============================================================================
\section{Bedingungsanalyse}

\subsection{Lerngruppe}
\begin{tabular}{ll}
    Kursgröße: & \lerngruppengroesse{} SuS (\maennlich{} m / \weiblich{} w / \divers{} d) \\
    Jahrgangsstufe: & \jahrgangsstufe \\
    Sprachniveau: & \gerniveau{} (GER) \\
    Fremdsprache: & 2. Fremdsprache \\
\end{tabular}

\subsection{Vorwissen}
Die SuS haben in der vorherigen Stunde (7a) die Grundlagen des Einkaufens kennengelernt und kennen bereits:
\begin{itemize}
    \item Zahlen 0-20
    \item Grundlegende Einkaufsdialoge
    \item Fragen nach dem Preis: ¿Cuánto cuesta?
\end{itemize}

\subsection{Differenzierung (intern)}

\textbf{WICHTIG:} Keine Sternchen auf Schülermaterial!

\begin{tabular}{|l|p{4cm}|p{4cm}|p{4cm}|}
\hline
& \textbf{[B] Basis} & \textbf{[S] Standard} & \textbf{[E] Erweiterung} \\
\hline
\textbf{Ziel} & Berufsreife & Mittlere Reife & Gymnasium \\
\hline
\textbf{Inhalt} & Zahlen erkennen, ablesen & Zahlen nennen, verstehen & Zahlen frei verwenden \\
\hline
\textbf{Hilfen} & Zahlentabelle, Aussprache & Teils mit Hilfe & Nur Impulse \\
\hline
\end{tabular}

\subsection{Besonderheiten}
\begin{itemize}
    \item \textbf{LRS:} 2 SuS (Nachteilsausgleich: Zeitzugabe)
    \item \textbf{DaZ:} 1 SuS (Wortschatzarbeit verstärken)
    \item \textbf{Leistungsstark:} 4 SuS (Zusatzaufgaben: Rechenaufgaben mit Preisen)
\end{itemize}

% ============================================================================
% 2. SACHANALYSE
% ============================================================================
\section{Sachanalyse}

\subsection{Sprachliche Strukturen}

\subsubsection{Zahlen 21-29}
Die Zahlen 21-29 werden als ein Wort geschrieben:
\begin{itemize}
    \item \textbf{veintiuno} (21) -- Achtung: vor mask. Subst. \textit{veintiún}
    \item \textbf{veintidós} (22) -- mit Akzent
    \item \textbf{veintitrés} (23) -- mit Akzent
    \item \textbf{veinticuatro} (24)
    \item \textbf{veinticinco} (25)
    \item \textbf{veintiséis} (26) -- mit Akzent
    \item \textbf{veintisiete} (27)
    \item \textbf{veintiocho} (28)
    \item \textbf{veintinueve} (29)
\end{itemize}

\subsubsection{Zehner 30-100}
\begin{itemize}
    \item \textbf{treinta} (30)
    \item \textbf{cuarenta} (40)
    \item \textbf{cincuenta} (50)
    \item \textbf{sesenta} (60)
    \item \textbf{setenta} (70)
    \item \textbf{ochenta} (80)
    \item \textbf{noventa} (90)
    \item \textbf{cien} (100) -- nur alleinstehend, sonst \textit{ciento}
\end{itemize}

\subsubsection{Kombinationen 31-99}
Zehner + \textbf{y} + Einer: \textit{treinta y uno, cuarenta y cinco, noventa y nueve}

\subsubsection{Preise}
\begin{itemize}
    \item \textit{¿Cuánto cuesta?} -- Wieviel kostet es?
    \item \textit{Cuesta X euros.} -- Es kostet X Euro.
    \item \textit{Cuesta X euros y Y céntimos.} -- Es kostet X Euro und Y Cent.
\end{itemize}

\subsection{Kontrastive Betrachtung (Deutsch $\leftrightarrow$ Spanisch)}
\begin{itemize}
    \item \textbf{Positive Transfers:} Euro als Währung identisch
    \item \textbf{Unterschiede:} Deutsche Reihenfolge bei Zahlen (einundzwanzig vs. veintiuno)
    \item \textbf{Stolpersteine:} Akzente bei 22, 23, 26; Schreibung als ein Wort (21-29)
\end{itemize}

\subsection{Kulturelle Aspekte}
\begin{itemize}
    \item Preise in Spanien: Dezimaltrennzeichen ist Komma (wie in Deutschland)
    \item Typische Preise auf Märkten und in Geschäften
\end{itemize}

% ============================================================================
% 3. DIDAKTISCHE ANALYSE
% ============================================================================
\section{Didaktische Analyse}

\subsection{Stellung in der Sequenz}
Dies ist die \textbf{7b. Doppelstunde} der Sequenz ``\reihenthema'' (\reihenstunden{} Doppelstunden gesamt).

Die SuS haben in 7a bereits Einkaufsdialoge kennengelernt und können nach Preisen fragen. Nun erweitern sie ihren Zahlenraum von 20 auf 100, um realistische Preise verstehen und nennen zu können.

\subsection{Kompetenzziele (Rahmenplan MV)}
\begin{itemize}
    \item \textbf{Kommunikativ:} Preise verstehen und nennen können
    \item \textbf{Sprachlich:} Zahlen 21-100 korrekt aussprechen und schreiben
    \item \textbf{Interkulturell:} Preise in spanischsprachigen Ländern
    \item \textbf{Methodisch:} Systematisches Lernen von Zahlensystemen
\end{itemize}

\subsection{Legitimation}
\textbf{Rahmenplan Spanisch MV (Kl. 7-10):}
\begin{itemize}
    \item Kompetenzbereich: Verfügung über sprachliche Mittel (Wortschatz)
    \item Standard: Zahlen und Preise im Alltag verstehen und verwenden
\end{itemize}

% ============================================================================
% 4. STUNDENZIELE
% ============================================================================
\section{Stundenziele}

\subsection{Grobziel}
Die SuS erweitern ihre \textbf{sprachliche Kompetenz} im Bereich Wortschatz, indem sie die Zahlen 21-100 verstehen, aussprechen und im Kontext von Preisangaben anwenden.

\subsection{Feinziele (operationalisiert)}
\begin{tabular}{|c|p{8cm}|c|c|}
\hline
\textbf{Nr.} & \textbf{Die SuS können...} & \textbf{AFB} & \textbf{Diff.} \\
\hline
FZ1 & die Zahlen 21-100 erkennen und zuordnen & I & alle \\
\hline
FZ2 & die Zahlen 21-100 korrekt aussprechen & I & alle \\
\hline
FZ3 & die Systematik der Zahlenbildung erklären & II & [S]/[E] \\
\hline
FZ4 & Preise verstehen und selbst nennen & II & alle \\
\hline
FZ5 & einen Einkaufsdialog mit Preisen führen & II/III & [S]/[E] \\
\hline
\end{tabular}

% ============================================================================
% 5. METHODISCHE ANALYSE
% ============================================================================
\section{Methodische Analyse}

\subsection{Phasierung (Doppelstunde = 2 × 45 Min.)}
\begin{enumerate}
    \item \textbf{1. Stunde:}
    \begin{itemize}
        \item Einstieg: Wiederholung Zahlen 0-20 (Blitzrunde)
        \item Input: Präsentation Zahlen 21-29 (zusammengeschrieben)
        \item Übung 1: Chorisches Sprechen
        \item Input: Zehner 30-100
        \item Übung 2: Zahlen-Bingo
    \end{itemize}
    \item \textbf{2. Stunde:}
    \begin{itemize}
        \item Wiederholung: Zahlen-Quiz
        \item Input: Preise mit Euro/Céntimos
        \item Anwendung: Einkaufsdialoge (Partnerarbeit)
        \item Sicherung: Arbeitsblatt
        \item Abschluss: Lernpfad-Hinweis
    \end{itemize}
\end{enumerate}

\subsection{Medien}
\begin{itemize}
    \item Tafel/Whiteboard
    \item Lehrwerk: \lehrwerk, \lehrwerkseite
    \item Arbeitsblatt: AB-07b.pdf
    \item Lernpfad: LP-07b.html (optional, digital)
    \item Bingo-Karten (vorbereitet)
    \item Preisschilder (Bildkarten)
\end{itemize}

% ============================================================================
% 6. VERLAUFSPLANUNG
% ============================================================================
\section{Verlaufsplanung}

\textbf{1. Stunde (45 Min.)}

\begin{longtable}{|p{1cm}|p{1.5cm}|p{4cm}|p{3cm}|p{2cm}|p{2cm}|}
\hline
\textbf{Zeit} & \textbf{Phase} & \textbf{Lehrerhandeln} & \textbf{Schülerhandeln} & \textbf{SF} & \textbf{Medien} \\
\hline
\endhead

0--5 & Einstieg & Begrüßung, Blitzrunde: Zahlen 0-20 abfragen & Zahlen nennen & Plenum & -- \\
\hline

5--15 & Input 1 & Tafelanschrieb: 21-29, Betonung der Schreibung als ein Wort & Zuhören, Notieren & Plenum & Tafel \\
\hline

15--22 & Übung 1 & Chorisches Sprechen anleiten: 21-29 & Nachsprechen, Rhythmus & Plenum & -- \\
\hline

22--30 & Input 2 & Zehner 30-100 einführen, Systematik zeigen & Zuhören, Muster erkennen & Plenum & Tafel \\
\hline

30--40 & Übung 2 & Zahlen-Bingo anleiten, Zahlen vorlesen & Zahlen ankreuzen & EA & Bingo-Karten \\
\hline

40--45 & Sicherung & Gewinner küren, schwierige Zahlen wiederholen & Kontrollieren & Plenum & -- \\
\hline
\end{longtable}

\vspace{1em}

\textbf{2. Stunde (45 Min.)}

\begin{longtable}{|p{1cm}|p{1.5cm}|p{4cm}|p{3cm}|p{2cm}|p{2cm}|}
\hline
\textbf{Zeit} & \textbf{Phase} & \textbf{Lehrerhandeln} & \textbf{Schülerhandeln} & \textbf{SF} & \textbf{Medien} \\
\hline
\endhead

0--5 & Wieder\-holung & Quiz: ``¿Qué número es?'' (Zahlen zeigen) & Zahlen nennen & Plenum & Karten \\
\hline

5--12 & Input 3 & Preise einführen: ``Cuesta X euros (y Y céntimos)'' & Zuhören, Beispiele & Plenum & Preisschilder \\
\hline

12--25 & Anwen\-dung & Partnerarbeit: Einkaufsdialoge mit Preisschildern & Dialoge üben & PA & Preisschilder \\
\hline

25--35 & Festigung & AB-07b ausgeben, bearbeiten lassen & AB bearbeiten & EA & AB \\
\hline

35--40 & Sicherung & Lösungen vergleichen, Fragen klären & Vergleichen, Fragen & Plenum & Tafel \\
\hline

40--45 & Abschluss & Zusammenfassung, HA, LP-07b empfehlen & Notieren & Plenum & -- \\
\hline
\end{longtable}

\textbf{Legende:} EA = Einzelarbeit, PA = Partnerarbeit, GA = Gruppenarbeit, SF = Sozialform

% ============================================================================
% 7. TAFELANSCHRIEB
% ============================================================================
\section{Tafelanschrieb}

\begin{infobox}[Tafelanschrieb]
\textbf{Los números 21-100}

\vspace{1em}

\textbf{21-29:} \textit{ein Wort!}

\begin{tabular}{llll}
21 veintiuno & 22 veintidós & 23 veintitrés & 24 veinticuatro \\
25 veinticinco & 26 veintiséis & 27 veintisiete & 28 veintiocho \\
29 veintinueve & & & \\
\end{tabular}

\vspace{1em}

\textbf{Die Zehner:}

\begin{tabular}{llll}
30 treinta & 40 cuarenta & 50 cincuenta & 60 sesenta \\
70 setenta & 80 ochenta & 90 noventa & 100 cien \\
\end{tabular}

\vspace{1em}

\textbf{Kombinationen:} Zehner + \textbf{y} + Einer

31 = treinta y uno \quad 45 = cuarenta y cinco \quad 99 = noventa y nueve

\vspace{1em}

\textbf{Los precios:}

¿Cuánto cuesta? -- Wieviel kostet es?

Cuesta \underline{\hspace{1cm}} euros (y \underline{\hspace{1cm}} céntimos).
\end{infobox}

% ============================================================================
% 8. ERWARTETE SCHWIERIGKEITEN
% ============================================================================
\section{Erwartete Schwierigkeiten}

\begin{tabular}{|p{5cm}|p{8cm}|}
\hline
\textbf{Schwierigkeit} & \textbf{Reaktion} \\
\hline
Schreibung 21-29 (ein Wort) & Visualisierung an Tafel, Eselsbrücke: ``veinti-'' + Einer \\
\hline
Akzente bei 22, 23, 26 & Betonung übertrieben sprechen, markieren \\
\hline
Verwechslung 40/50 (cuarenta/cincuenta) & Kontrastive Übung, Bewegungsspiel \\
\hline
``y'' vergessen bei 31-99 & Immer laut mitsprechen: ``treinta Y uno'' \\
\hline
Céntimos-Aussprache & Chorisches Sprechen, th-Laut üben \\
\hline
Zeitproblem & Puffer: Bingo kürzen, Dialog-Präsentation freiwillig \\
\hline
\end{tabular}

% ============================================================================
% 9. HAUSAUFGABE
% ============================================================================
\section{Hausaufgabe}

\begin{itemize}
    \item Lehrwerk S. 49, Aufgabe 3 (Zahlen schreiben)
    \item Vokabeln lernen: Zahlen 21-100, Preisausdrücke
    \item Optional: Lernpfad LP-07b.html bearbeiten (Tablet/PC)
\end{itemize}

% ============================================================================
% 10. ANLAGEN
% ============================================================================
\section{Anlagen}

\begin{enumerate}
    \item Arbeitsblatt (AB-07b.pdf)
    \item Musterlösung (ML-07b.pdf)
    \item Lehrerhinweise (LH-07b.pdf)
    \item Lernpfad (LP-07b.html)
    \item Bingo-Karten (25 verschiedene, Zahlen 21-100)
    \item Preisschilder (Bildkarten mit Produkten und Preisen)
\end{enumerate}

% ============================================================================
% 11. DIDAKTIK-SCORE
% ============================================================================
\newpage
\section{Didaktik-Score}

\begin{scorebox}
\textbf{Bewertung der didaktischen Qualität nach 10 Dimensionen}\\
\textbf{Ziel-Score: $\geq$ 9.0 (Premium-Qualität)}

\vspace{1em}
\begin{tabular}{|c|l|c|c|p{4.5cm}|}
\hline
\textbf{\#} & \textbf{Dimension} & \textbf{Gew.} & \textbf{Score} & \textbf{Notiz} \\
\hline
1 & Differenzierung & 15\% & 9/10 & [B]/[S]/[E] klar definiert \\
\hline
2 & Taxonomie (AFB) & 10\% & 9/10 & AFB I-III ausgewogen \\
\hline
3 & Lernziel-Alignment & 10\% & 10/10 & RP MV erfüllt \\
\hline
4 & Aktivierung & 10\% & 9/10 & Bingo, Dialoge, chorisch \\
\hline
5 & Scaffolding & 10\% & 9/10 & Systematischer Aufbau \\
\hline
6 & Feedback-Qualität & 10\% & 9/10 & Sofort bei Übungen \\
\hline
7 & Sprachliche Klarheit & 10\% & 10/10 & Altersgerecht \\
\hline
8 & Motivation \& Relevanz & 10\% & 9/10 & Alltagsbezug: Einkaufen \\
\hline
9 & Inklusion (UDL) & 10\% & 9/10 & Visuell + auditiv + handelnd \\
\hline
10 & Praktikabilität & 5\% & 9/10 & Realistisch für 90 Min. \\
\hline
\multicolumn{3}{|r|}{\textbf{GESAMT:}} & \textbf{9.2/10} & \\
\hline
\end{tabular}

\vspace{1em}
\textbf{Status:} \textcolor{successgreen}{\textbf{FREIGABE}} (Premium-Qualität erreicht)
\end{scorebox}

\end{document}
